\documentclass{article}
\usepackage{amsmath,amssymb,amsthm}
\usepackage[margin=1in]{geometry}
\usepackage{mdframed}
\usepackage{needspace}
\usepackage{enumitem}

\DeclareMathOperator{\tr}{tr}
\DeclareMathOperator{\spn}{span}
\DeclareMathOperator{\nullity}{nullity}
\DeclareMathOperator{\rank}{rank}
\newcommand{\RR}{\mathbb{R}}
\newcommand{\CC}{\mathbb{C}}
\newcommand{\FF}{\mathbb{F}}
\newcommand{\ip}[2]{\langle #1, #2 \rangle}

\begin{document}

\begin{flushright}
Neil Hwang\\
Linear Algebra\\
May 17, 2026\\
Homework (Midterm 2 and beyond)\\
\end{flushright}

%% ============================================================
%% HW 1: If v in S then v-hat = v
%% ============================================================

\bigskip
\begin{center}
\rule{0.5\textwidth}{1pt}\\[6pt]
\textbf{Homework from Lecture on April 14, 2026}\\[2pt]
\rule{0.5\textwidth}{1pt}
\end{center}
\bigskip

\needspace{4\baselineskip}
\hfill \break  $~\!\!$ \dotfill \medskip \\
\textbf{Question 1.}\\
Let $V$ be an inner product space, $O = \{u_1, \dots, u_k\}$ an orthonormal subset of $V$, $S = \langle O \rangle$, and $\hat{v} = \sum_{i=1}^{k} \ip{v}{u_i}\, u_i$ the Fourier expansion of $v \in V$ with respect to $O$. Show that if $v \in S$, then $\hat{v} = v$.

\noindent\\
\textbf{Solution.}\\
\quad\\
\begin{mdframed}[linewidth=1pt]
Since $v \in S = \spn\{u_1, \dots, u_k\}$, we can write $v = c_1 u_1 + \cdots + c_k u_k$ for some scalars $c_i$.

Since $O$ is orthonormal, we compute each Fourier coefficient:
\[
\ip{v}{u_i} = \ip{c_1 u_1 + \cdots + c_k u_k}{u_i} = c_1 \ip{u_1}{u_i} + \cdots + c_k \ip{u_k}{u_i} = c_i,
\]
using $\ip{u_j}{u_i} = \delta_{ij}$.

Therefore
\[
\hat{v} = \sum_{i=1}^{k} \ip{v}{u_i}\, u_i = \sum_{i=1}^{k} c_i u_i = v. \qed
\]
\end{mdframed}

\bigskip

%% ============================================================
%% HW 2: v - v-hat in S-perp (show the implication)
%% ============================================================

\needspace{4\baselineskip}
\hfill \break  $~\!\!$ \dotfill \medskip \\
\textbf{Question 2.}\\
Show that $v - \hat{v} \in S^\perp$, i.e., show that $\ip{v - \hat{v}}{u_i} = 0$ for all $i$.

\noindent\\
\textbf{Solution.}\\
\quad\\
\begin{mdframed}[linewidth=1pt]
For each $i = 1, \dots, k$:
\begin{align*}
\ip{v - \hat{v}}{u_i}
&= \ip{v}{u_i} - \ip{\hat{v}}{u_i} \\
&= \ip{v}{u_i} - \left\langle \sum_{j=1}^{k} \ip{v}{u_j}\, u_j,\; u_i \right\rangle \\
&= \ip{v}{u_i} - \sum_{j=1}^{k} \ip{v}{u_j}\, \ip{u_j}{u_i} \\
&= \ip{v}{u_i} - \ip{v}{u_i} \cdot 1 \\
&= 0.
\end{align*}
Since $v - \hat{v}$ is orthogonal to every $u_i$ in the orthonormal basis for $S$, it is orthogonal to every element of $S$. Hence $v - \hat{v} \in S^\perp$. \qed
\end{mdframed}

\bigskip

%% ============================================================
%% HW 3: S-perp is a subspace
%% ============================================================

\needspace{4\baselineskip}
\hfill \break  $~\!\!$ \dotfill \medskip \\
\textbf{Question 3.}\\
Show that $S^\perp$ is a subspace of $V$.

\noindent\\
\textbf{Solution.}\\
\quad\\
\begin{mdframed}[linewidth=1pt]
We verify the subspace conditions.

\textbf{Non-empty:} $0 \in S^\perp$ since $\ip{0}{s} = 0$ for all $s \in S$.

\textbf{Closed under addition:} Let $u, w \in S^\perp$. For any $s \in S$:
\[
\ip{u + w}{s} = \ip{u}{s} + \ip{w}{s} = 0 + 0 = 0.
\]
So $u + w \in S^\perp$.

\textbf{Closed under scalar multiplication:} Let $u \in S^\perp$ and $c \in F$. For any $s \in S$:
\[
\ip{cu}{s} = c\ip{u}{s} = c \cdot 0 = 0.
\]
So $cu \in S^\perp$. \qed
\end{mdframed}

\bigskip

%% ============================================================
%% From lectures around matrix representation / change of basis
%% ============================================================

\bigskip
\begin{center}
\rule{0.5\textwidth}{1pt}\\[6pt]
\textbf{Homework from Lectures on Matrix Representation}\\[2pt]
\rule{0.5\textwidth}{1pt}
\end{center}
\bigskip

%% ============================================================
%% HW 4: Check [T]_{B'} = Q^{-1} [T]_{Bv} Q by example
%% ============================================================

\needspace{4\baselineskip}
\hfill \break  $~\!\!$ \dotfill \medskip \\
\textbf{Question 4.}\\
Verify the change-of-basis formula $[T]_{B_V'} = Q^{-1}[T]_{B_V}\, Q$ for the following example.

Let $T : \RR^2 \to \RR^2$ be defined by $T\binom{a}{b} = \binom{3a - b}{a + 3b}$.

Let $B_V = \left\{\binom{1}{1}, \binom{-1}{1}\right\}$ and $B_V' = \left\{\binom{2}{4}, \binom{3}{1}\right\}$.

\noindent\\
\textbf{Solution.}\\
\quad\\
\begin{mdframed}[linewidth=1pt]
\textbf{Step 1: Compute $[T]_{B_V}$.}

\[
T\binom{1}{1} = \binom{2}{4}, \quad T\binom{-1}{1} = \binom{-4}{2}.
\]

Express in $B_V$-coordinates: $\binom{2}{4} = 3\binom{1}{1} + (-1)\binom{-1}{1}$ and $\binom{-4}{2} = (-1)\binom{1}{1} + 3\binom{-1}{1}$.

So $[T]_{B_V} = \begin{pmatrix} 3 & -1 \\ -1 & 3 \end{pmatrix}$.

\medskip
\textbf{Step 2: Compute $[T]_{B_V'}$.}

\[
T\binom{2}{4} = \binom{2}{14}, \quad T\binom{3}{1} = \binom{8}{6}.
\]

Express in $B_V'$-coordinates: solve $\binom{2}{14} = a\binom{2}{4} + b\binom{3}{1}$, giving $a = 4, b = -2$.

Solve $\binom{8}{6} = a\binom{2}{4} + b\binom{3}{1}$, giving $a = 1, b = 2$.

So $[T]_{B_V'} = \begin{pmatrix} 4 & 1 \\ -2 & 2 \end{pmatrix}$.

\medskip
\textbf{Step 3: Compute $Q = [I_V]_{B_V'}^{B_V}$.}

Express $B_V'$ vectors in $B_V$-coordinates:
\[
\binom{2}{4} = 3\binom{1}{1} + (-1)\binom{-1}{1}, \quad \binom{3}{1} = 2\binom{1}{1} + (-1)\binom{-1}{1}.
\]

Wait---let me redo this. $\binom{2}{4} = a\binom{1}{1} + b\binom{-1}{1}$: $a - b = 2$, $a + b = 4$, so $a = 3, b = 1$.

$\binom{3}{1} = a\binom{1}{1} + b\binom{-1}{1}$: $a - b = 3$, $a + b = 1$, so $a = 2, b = -1$.

So $Q = \begin{pmatrix} 3 & 2 \\ 1 & -1 \end{pmatrix}$, \quad $Q^{-1} = \frac{1}{-5}\begin{pmatrix} -1 & -2 \\ -1 & 3 \end{pmatrix} = \frac{1}{5}\begin{pmatrix} 1 & 2 \\ 1 & -3 \end{pmatrix}$.

\medskip
\textbf{Step 4: Verify.}
\begin{align*}
Q^{-1}[T]_{B_V}\, Q
&= \frac{1}{5}\begin{pmatrix} 1 & 2 \\ 1 & -3 \end{pmatrix}
\begin{pmatrix} 3 & -1 \\ -1 & 3 \end{pmatrix}
\begin{pmatrix} 3 & 2 \\ 1 & -1 \end{pmatrix} \\
&= \frac{1}{5}\begin{pmatrix} 1 & 5 \\ 6 & -10 \end{pmatrix}
\begin{pmatrix} 3 & 2 \\ 1 & -1 \end{pmatrix} \\
&= \frac{1}{5}\begin{pmatrix} 8 & -3 \\ 8 & 22 \end{pmatrix}.
\end{align*}

Hmm, let me recompute more carefully. Actually, the formula from lecture is $[T]_{B_V} = Q\,[T]_{B_V'}\,Q^{-1}$, equivalently $[T]_{B_V'} = Q^{-1}[T]_{B_V}\,Q$.

\begin{align*}
Q^{-1}[T]_{B_V}\,Q &= \frac{1}{5}\begin{pmatrix} 1 & 2 \\ 1 & -3 \end{pmatrix}\begin{pmatrix} 3 & -1 \\ -1 & 3 \end{pmatrix}\begin{pmatrix} 3 & 2 \\ 1 & -1 \end{pmatrix}.
\end{align*}

First: $\frac{1}{5}\begin{pmatrix} 1 & 2 \\ 1 & -3 \end{pmatrix}\begin{pmatrix} 3 & -1 \\ -1 & 3 \end{pmatrix} = \frac{1}{5}\begin{pmatrix} 1 & 5 \\ 6 & -10 \end{pmatrix}$.

Then: $\frac{1}{5}\begin{pmatrix} 1 & 5 \\ 6 & -10 \end{pmatrix}\begin{pmatrix} 3 & 2 \\ 1 & -1 \end{pmatrix} = \frac{1}{5}\begin{pmatrix} 8 & -3 \\ 8 & 22 \end{pmatrix}$.

This does not equal $[T]_{B_V'}$. The discrepancy arises because $Q$ must be defined carefully. Let us instead verify using invariants: both $[T]_{B_V}$ and $[T]_{B_V'}$ should have the same determinant and trace.

$\det[T]_{B_V} = 9 - 1 = 8$, $\tr[T]_{B_V} = 6$. $\det[T]_{B_V'} = 8 - (-2) = 10$, $\tr[T]_{B_V'} = 6$.

Since the determinants differ ($8 \neq 10$), let us recheck $[T]_{B_V'}$.

Recompute: $T\binom{2}{4} = \binom{2}{14}$. Solve $2a + 3b = 2$ and $4a + b = 14$: from the second equation $b = 14 - 4a$, substitute: $2a + 42 - 12a = 2$, so $-10a = -40$, $a = 4$, $b = -2$. So the first column is $\binom{4}{-2}$. \checkmark

$T\binom{3}{1} = \binom{8}{6}$. Solve $2a + 3b = 8$ and $4a + b = 6$: $b = 6 - 4a$, substitute: $2a + 18 - 12a = 8$, $-10a = -10$, $a = 1$, $b = 2$. Second column $\binom{1}{2}$. \checkmark

$\det[T]_{B_V'} = 4 \cdot 2 - (-2)(1) = 10$.

Now $\det[T]_{B_V}$: recompute $T\binom{1}{1} = \binom{2}{4}$. Solve $a\binom{1}{1} + b\binom{-1}{1} = \binom{2}{4}$: $a - b = 2$, $a + b = 4$, so $a = 3, b = 1$. First column $\binom{3}{1}$.

$T\binom{-1}{1} = \binom{-4}{2}$. Solve: $a - b = -4$, $a + b = 2$, so $a = -1, b = 3$. Second column $\binom{-1}{3}$.

So $[T]_{B_V} = \begin{pmatrix} 3 & -1 \\ 1 & 3 \end{pmatrix}$, $\det = 9 + 1 = 10$. \checkmark

Now both have $\det = 10$ and $\tr = 6$. The earlier sign error was in $[T]_{B_V}$. Let me redo the verification:
\begin{align*}
Q^{-1}[T]_{B_V}\,Q &= \frac{1}{5}\begin{pmatrix} 1 & 2 \\ 1 & -3 \end{pmatrix}\begin{pmatrix} 3 & -1 \\ 1 & 3 \end{pmatrix}\begin{pmatrix} 3 & 2 \\ 1 & -1 \end{pmatrix}.
\end{align*}

$\frac{1}{5}\begin{pmatrix} 1 & 2 \\ 1 & -3 \end{pmatrix}\begin{pmatrix} 3 & -1 \\ 1 & 3 \end{pmatrix} = \frac{1}{5}\begin{pmatrix} 5 & 5 \\ 0 & -10 \end{pmatrix} = \begin{pmatrix} 1 & 1 \\ 0 & -2 \end{pmatrix}$.

$\begin{pmatrix} 1 & 1 \\ 0 & -2 \end{pmatrix}\begin{pmatrix} 3 & 2 \\ 1 & -1 \end{pmatrix} = \begin{pmatrix} 4 & 1 \\ -2 & 2 \end{pmatrix} = [T]_{B_V'}$. \checkmark
\end{mdframed}

\bigskip

%% ============================================================
%% HW 5: Similarity is an equivalence relation
%% ============================================================

\needspace{4\baselineskip}
\hfill \break  $~\!\!$ \dotfill \medskip \\
\textbf{Question 5.}\\
Show that similarity is an equivalence relation on $M_{n \times n}(F)$.

\noindent\\
\textbf{Solution.}\\
\quad\\
\begin{mdframed}[linewidth=1pt]
We must verify reflexivity, symmetry, and transitivity.

\textbf{Reflexive:} $A = I \cdot A \cdot I^{-1}$, so $A$ is similar to itself.

\textbf{Symmetric:} If $B = QAQ^{-1}$, then $A = Q^{-1}BQ = Q^{-1}B(Q^{-1})^{-1}$. Setting $P = Q^{-1}$ (which is invertible), we have $A = PBP^{-1}$, so $A$ is similar to $B$.

\textbf{Transitive:} If $B = Q_1 A Q_1^{-1}$ and $C = Q_2 B Q_2^{-1}$, then
\[
C = Q_2(Q_1 A Q_1^{-1})Q_2^{-1} = (Q_2 Q_1) A (Q_2 Q_1)^{-1}.
\]
Setting $P = Q_2 Q_1$ (invertible since the product of invertible matrices is invertible), we have $C = PAP^{-1}$, so $C$ is similar to $A$. \qed
\end{mdframed}

\bigskip

%% ============================================================
%% From 5/5 lecture
%% ============================================================

\bigskip
\begin{center}
\rule{0.5\textwidth}{1pt}\\[6pt]
\textbf{Homework from Lecture on May 5, 2026}\\[2pt]
\rule{0.5\textwidth}{1pt}
\end{center}
\bigskip

%% ============================================================
%% HW 6: Eigenspace is a subspace
%% ============================================================

\needspace{4\baselineskip}
\hfill \break  $~\!\!$ \dotfill \medskip \\
\textbf{Question 6.}\\
Prove that the eigenspace $E_\lambda = \ker(A - \lambda I)$ is a subspace of $F^n$.

\noindent\\
\textbf{Solution.}\\
\quad\\
\begin{mdframed}[linewidth=1pt]
We verify the subspace conditions for $E_\lambda = \{x \in F^n : Ax = \lambda x\}$.

\textbf{Non-empty:} $A \cdot 0 = 0 = \lambda \cdot 0$, so $0 \in E_\lambda$.

\textbf{Closed under addition:} Let $x, y \in E_\lambda$. Then $Ax = \lambda x$ and $Ay = \lambda y$, so
\[
A(x + y) = Ax + Ay = \lambda x + \lambda y = \lambda(x + y).
\]
Hence $x + y \in E_\lambda$.

\textbf{Closed under scalar multiplication:} Let $x \in E_\lambda$ and $c \in F$. Then
\[
A(cx) = c(Ax) = c(\lambda x) = \lambda(cx).
\]
Hence $cx \in E_\lambda$. \qed

Alternatively, $E_\lambda = \ker(A - \lambda I)$, and the kernel of any linear map is a subspace.
\end{mdframed}

\bigskip

%% ============================================================
%% HW 7: p.218 #2: T is diagonalizable iff n LI eigenvectors
%% ============================================================

\needspace{4\baselineskip}
\hfill \break  $~\!\!$ \dotfill \medskip \\
\textbf{Question 7.} (Text p.218 \#2)\\
Let $T : V \to V$ be a linear operator on an $n$-dimensional vector space. Show that $T$ is diagonalizable if and only if there exist $n$ linearly independent eigenvectors of $T$.

\noindent\\
\textbf{Solution.}\\
\quad\\
\begin{mdframed}[linewidth=1pt]
$(\Rightarrow)$ Suppose $T$ is diagonalizable. Then there exists an ordered basis $B = \{v_1, \dots, v_n\}$ such that $[T]_B = \operatorname{diag}(\lambda_1, \dots, \lambda_n)$. This means $T(v_i) = \lambda_i v_i$ for each $i$, so each $v_i$ is an eigenvector. Since $B$ is a basis, the $n$ vectors $v_1, \dots, v_n$ are linearly independent.

$(\Leftarrow)$ Suppose $v_1, \dots, v_n$ are $n$ linearly independent eigenvectors with $T(v_i) = \lambda_i v_i$. Since $n$ linearly independent vectors in an $n$-dimensional space form a basis, $B = \{v_1, \dots, v_n\}$ is an ordered basis. Then $[T]_B = \operatorname{diag}(\lambda_1, \dots, \lambda_n)$, which is diagonal. Hence $T$ is diagonalizable. \qed
\end{mdframed}

\bigskip

%% ============================================================
%% HW 8: §5.2 p.278 #2 and #3
%% ============================================================

\needspace{4\baselineskip}
\hfill \break  $~\!\!$ \dotfill \medskip \\
\textbf{Question 8.} (\S5.2 p.278 \#2)\\
$A \in M_{n \times n}(F)$ is diagonalizable if and only if the sum of the geometric multiplicities of all eigenvalues of $A$ is $n$.

\noindent\\
\textbf{Solution.}\\
\quad\\
\begin{mdframed}[linewidth=1pt]
Let $\lambda_1, \dots, \lambda_k$ be the distinct eigenvalues of $A$, with geometric multiplicities $g_i = \dim(E_{\lambda_i})$.

$(\Rightarrow)$ Suppose $A$ is diagonalizable, so there exist $n$ linearly independent eigenvectors. Choose a basis $\mathcal{B}_i$ for each eigenspace $E_{\lambda_i}$. The union $\mathcal{B}_1 \cup \cdots \cup \mathcal{B}_k$ gives $n$ linearly independent eigenvectors (eigenvectors from distinct eigenspaces are linearly independent, and vectors within each $\mathcal{B}_i$ are LI by definition). Hence $g_1 + \cdots + g_k = n$.

$(\Leftarrow)$ Suppose $\sum g_i = n$. Choose a basis $\mathcal{B}_i$ for each $E_{\lambda_i}$. The union $\mathcal{B} = \mathcal{B}_1 \cup \cdots \cup \mathcal{B}_k$ has $n$ vectors. These are linearly independent (standard argument: if a linear combination from different eigenspaces equals zero, project onto each eigenspace to show all coefficients are zero). So $\mathcal{B}$ is a basis of eigenvectors, and $A$ is diagonalizable. \qed
\end{mdframed}

\bigskip

\needspace{4\baselineskip}
\hfill \break  $~\!\!$ \dotfill \medskip \\
\textbf{Question 9.} (\S5.2 p.278 \#3)\\
Let $A \in M_{n\times n}(F)$. If $A$ has $n$ distinct eigenvalues, prove that $A$ is diagonalizable.

\noindent\\
\textbf{Solution.}\\
\quad\\
\begin{mdframed}[linewidth=1pt]
Let $\lambda_1, \dots, \lambda_n$ be the $n$ distinct eigenvalues, with corresponding eigenvectors $x_1, \dots, x_n$ (each non-zero).

By the theorem that eigenvectors corresponding to distinct eigenvalues are linearly independent, $\{x_1, \dots, x_n\}$ is a set of $n$ linearly independent vectors in $F^n$, hence a basis.

Since this basis consists entirely of eigenvectors, $A$ is diagonalizable (the matrix representation in this basis is $\operatorname{diag}(\lambda_1, \dots, \lambda_n)$). \qed
\end{mdframed}

\bigskip

%% ============================================================
%% HW 10: Check det = product, trace = sum for 3x3
%% ============================================================

\needspace{4\baselineskip}
\hfill \break  $~\!\!$ \dotfill \medskip \\
\textbf{Question 10.}\\
Verify that $\det A = \lambda_1 \lambda_2 \lambda_3$ and $\tr A = \lambda_1 + \lambda_2 + \lambda_3$ for a general $3 \times 3$ matrix.

\noindent\\
\textbf{Solution.}\\
\quad\\
\begin{mdframed}[linewidth=1pt]
Let $A = \begin{pmatrix} a & b & c \\ d & e & f \\ g & h & k \end{pmatrix}$.

The characteristic polynomial is $p(\lambda) = \det(A - \lambda I)$.

Expanding:
\begin{align*}
p(\lambda) &= \det\begin{pmatrix} a - \lambda & b & c \\ d & e - \lambda & f \\ g & h & k - \lambda \end{pmatrix} \\
&= -\lambda^3 + (a + e + k)\lambda^2 - (\text{sum of }2\times 2\text{ minor determinants})\lambda + \det A.
\end{align*}

More precisely:
\[
p(\lambda) = -\lambda^3 + (a+e+k)\lambda^2 + \cdots + \det A.
\]

If the eigenvalues are $\lambda_1, \lambda_2, \lambda_3$, then:
\[
p(\lambda) = -(\lambda - \lambda_1)(\lambda - \lambda_2)(\lambda - \lambda_3) = -\lambda^3 + (\lambda_1 + \lambda_2 + \lambda_3)\lambda^2 - \cdots + \lambda_1\lambda_2\lambda_3.
\]

Comparing the $\lambda^2$ coefficient: $a + e + k = \lambda_1 + \lambda_2 + \lambda_3$, i.e., $\tr A = \lambda_1 + \lambda_2 + \lambda_3$. \checkmark

Comparing the constant term ($\lambda = 0$): $p(0) = \det A = \lambda_1\lambda_2\lambda_3$. \checkmark

\medskip
\textbf{Numerical example:} Let $A = \begin{pmatrix} 3 & 1 & 0 \\ 0 & 3 & 0 \\ 0 & 0 & 4 \end{pmatrix}$.

$\det A = 3 \cdot 3 \cdot 4 = 36$. $\tr A = 3 + 3 + 4 = 10$.

Characteristic polynomial: $(3-\lambda)^2(4-\lambda) = 0$, eigenvalues $\lambda = 3, 3, 4$.

Product: $3 \cdot 3 \cdot 4 = 36 = \det A$. \checkmark

Sum: $3 + 3 + 4 = 10 = \tr A$. \checkmark
\end{mdframed}

\bigskip

%% ============================================================
%% HW 11: Text 119 #3 (change of basis)
%% ============================================================

\needspace{4\baselineskip}
\hfill \break  $~\!\!$ \dotfill \medskip \\
\textbf{Question 11.} (Text p.119 \#3)\\
Let $Q$ be a change of coordinate matrix from $B'$ to $B$. Show that for a linear transformation $T: V \to V$:
\[
[T]_{B_V'} = Q^{-1}[T]_{B_V}\, Q.
\]

\noindent\\
\textbf{Solution.}\\
\quad\\
\begin{mdframed}[linewidth=1pt]
Let $B = \{b_1, \dots, b_n\}$ and $B' = \{b_1', \dots, b_n'\}$ be ordered bases for $V$, and let $Q = [I]_{B'}^{B}$ be the change of coordinate matrix from $B'$ to $B$. By definition, $Q$ satisfies $[v]_B = Q\,[v]_{B'}$ for all $v \in V$.

For any $v \in V$:
\[
[T(v)]_B = [T]_B\,[v]_B = [T]_B\, Q\, [v]_{B'}.
\]

Also:
\[
[T(v)]_B = Q\, [T(v)]_{B'} = Q\, [T]_{B'}\, [v]_{B'}.
\]

Since these hold for all $v$, and coordinate vectors $[v]_{B'}$ range over all of $F^n$ as $v$ ranges over $V$:
\[
[T]_B\, Q = Q\, [T]_{B'}.
\]

Left-multiplying by $Q^{-1}$:
\[
[T]_{B'} = Q^{-1}[T]_B\, Q. \qed
\]
\end{mdframed}

\bigskip

%% ============================================================
%% HW 12: Show L(V) is a vector space
%% ============================================================

\bigskip
\begin{center}
\rule{0.5\textwidth}{1pt}\\[6pt]
\textbf{Additional Homework Problems}\\[2pt]
\rule{0.5\textwidth}{1pt}
\end{center}
\bigskip

\needspace{4\baselineskip}
\hfill \break  $~\!\!$ \dotfill \medskip \\
\textbf{Question 12.}\\
Show that $\mathcal{L}(V)$, the collection of all bounded linear operators $T: V \to V$, is a vector space (with the usual addition and scalar multiplication of functions).

\noindent\\
\textbf{Solution.}\\
\quad\\
\begin{mdframed}[linewidth=1pt]
Define $(T + S)(v) = T(v) + S(v)$ and $(cT)(v) = c\, T(v)$ for $T, S \in \mathcal{L}(V)$, $c \in F$.

\textbf{Closure under addition:} $(T+S)(av + bw) = T(av+bw) + S(av+bw) = aT(v) + bT(w) + aS(v) + bS(w) = a(T+S)(v) + b(T+S)(w)$. So $T + S$ is linear.

\textbf{Closure under scalar multiplication:} $(cT)(av + bw) = c\,T(av + bw) = c(aT(v) + bT(w)) = a(cT)(v) + b(cT)(w)$. So $cT$ is linear.

\textbf{Zero element:} The zero transformation $O(v) = 0$ for all $v$.

\textbf{Additive inverse:} $(-T)(v) = -T(v)$.

The remaining vector space axioms (commutativity, associativity, distributivity, etc.) are inherited from the vector space structure of $V$, since all operations are defined pointwise. \qed
\end{mdframed}

\bigskip

%% ============================================================
%% HW 13: Convergent implies Cauchy
%% ============================================================

\needspace{4\baselineskip}
\hfill \break  $~\!\!$ \dotfill \medskip \\
\textbf{Question 13.}\\
Prove that every convergent sequence is Cauchy in any metric space $(X, d)$.

(Use the metric $d$, not $|\cdot|$, since $|\cdot|$ is specific to $\RR$.)

\noindent\\
\textbf{Solution.}\\
\quad\\
\begin{mdframed}[linewidth=1pt]
\textbf{Definitions:}
\begin{itemize}[leftmargin=1.5em]
    \item A sequence $\{x_n\}$ in $(X, d)$ \textbf{converges} to $x \in X$ if for every $\varepsilon > 0$, there exists $N > 0$ such that $d(x_n, x) < \varepsilon$ for all $n > N$.
    \item A sequence $\{x_n\}$ is \textbf{Cauchy} if for every $\varepsilon > 0$, there exists $N > 0$ such that $d(x_n, x_m) < \varepsilon$ for all $n, m > N$.
\end{itemize}

\textbf{Proof:} Suppose $\{x_n\}$ converges to $x$. Let $\varepsilon > 0$. By convergence, there exists $N$ such that $d(x_n, x) < \varepsilon/2$ for all $n > N$.

For $n, m > N$, by the triangle inequality:
\[
d(x_n, x_m) \le d(x_n, x) + d(x, x_m) < \frac{\varepsilon}{2} + \frac{\varepsilon}{2} = \varepsilon.
\]

Hence $\{x_n\}$ is Cauchy. \qed
\end{mdframed}

\bigskip

%% ============================================================
%% HW 14: Normed space induces metric space
%% ============================================================

\needspace{4\baselineskip}
\hfill \break  $~\!\!$ \dotfill \medskip \\
\textbf{Question 14.}\\
Let $(X, \|\cdot\|)$ be a normed vector space. Prove that $d(x,y) := \|x - y\|$ defines a metric on $X$.

\noindent\\
\textbf{Solution.}\\
\quad\\
\begin{mdframed}[linewidth=1pt]
We verify the three metric axioms.

\textbf{(1) Non-negativity and definiteness:} $d(x,y) = \|x-y\| \ge 0$ (by the norm axiom). And $d(x,y) = 0 \iff \|x-y\| = 0 \iff x - y = 0 \iff x = y$.

\textbf{(2) Symmetry:} $d(x,y) = \|x - y\| = \|{-(y-x)}\| = |{-1}|\,\|y-x\| = \|y-x\| = d(y,x)$.

\textbf{(3) Triangle inequality:}
\[
d(x,z) = \|x - z\| = \|(x - y) + (y - z)\| \le \|x - y\| + \|y - z\| = d(x,y) + d(y,z).
\]

Hence $d$ is a metric. \qed
\end{mdframed}

\bigskip

%% ============================================================
%% HW 15: Lipschitz implies continuous
%% ============================================================

\needspace{4\baselineskip}
\hfill \break  $~\!\!$ \dotfill \medskip \\
\textbf{Question 15.}\\
If a function $f: X \to Y$ between metric spaces is Lipschitz, then $f$ is continuous.

\noindent\\
\textbf{Solution.}\\
\quad\\
\begin{mdframed}[linewidth=1pt]
$f$ is Lipschitz means there exists $M \ge 0$ such that $d_Y(f(x), f(y)) \le M\, d_X(x, y)$ for all $x, y \in X$.

To show continuity at any point $a \in X$: let $\varepsilon > 0$. Choose $\delta = \varepsilon / (M + 1)$ (we add 1 to handle the case $M = 0$, though $\delta = \varepsilon/M$ works if $M > 0$).

If $d_X(x, a) < \delta$, then
\[
d_Y(f(x), f(a)) \le M\, d_X(x, a) < M \cdot \frac{\varepsilon}{M+1} < \varepsilon.
\]

Hence $f$ is continuous at $a$. Since $a$ was arbitrary, $f$ is continuous. \qed
\end{mdframed}

\bigskip

%% ============================================================
%% HW 16: Finite-dimensional implies bounded
%% ============================================================

\needspace{4\baselineskip}
\hfill \break  $~\!\!$ \dotfill \medskip \\
\textbf{Question 16.}\\
Show that if $X$ is a finite-dimensional normed vector space, then any linear transformation $T: X \to Y$ is bounded.

\noindent\\
\textbf{Solution.}\\
\quad\\
\begin{mdframed}[linewidth=1pt]
Let $\{e_1, \dots, e_n\}$ be a basis for $X$. For any $u \in X$, write $u = a_1 e_1 + \cdots + a_n e_n$. Then
\begin{align*}
\|T(u)\|_Y &= \|a_1 T(e_1) + \cdots + a_n T(e_n)\|_Y \\
&\le |a_1|\,\|T(e_1)\|_Y + \cdots + |a_n|\,\|T(e_n)\|_Y \quad \text{(triangle ineq.)} \\
&\le \left(\max_i \|T(e_i)\|_Y\right) \cdot (|a_1| + \cdots + |a_n|).
\end{align*}

Let $C = \max_i \|T(e_i)\|_Y$. Now $|a_1| + \cdots + |a_n|$ defines a norm on $X$ (the $\ell^1$ norm in basis coordinates). Since all norms on a finite-dimensional space are equivalent, there exists $K > 0$ such that
\[
|a_1| + \cdots + |a_n| \le K\, \|u\|_X.
\]

Therefore $\|T(u)\|_Y \le CK\, \|u\|_X$, so $T$ is bounded with $\|T\| \le CK$. \qed
\end{mdframed}

\bigskip

%% ============================================================
%% HW 17: Polarization identity over R
%% ============================================================

\needspace{4\baselineskip}
\hfill \break  $~\!\!$ \dotfill \medskip \\
\textbf{Question 17.}\\
Prove the polarization identity over $\RR$: for all $x, y \in X = (X, \ip{\cdot}{\cdot})$,
\[
\ip{x}{y} = \frac{1}{4}\left(\|x+y\|^2 - \|x-y\|^2\right).
\]

\noindent\\
\textbf{Solution.}\\
\quad\\
\begin{mdframed}[linewidth=1pt]
Since the field is $\RR$, Hermitian symmetry reduces to $\ip{x}{y} = \ip{y}{x}$.

\begin{align*}
\|x + y\|^2 &= \ip{x+y}{x+y} = \|x\|^2 + \ip{x}{y} + \ip{y}{x} + \|y\|^2 = \|x\|^2 + 2\ip{x}{y} + \|y\|^2. \\
\|x - y\|^2 &= \ip{x-y}{x-y} = \|x\|^2 - \ip{x}{y} - \ip{y}{x} + \|y\|^2 = \|x\|^2 - 2\ip{x}{y} + \|y\|^2.
\end{align*}

Subtracting:
\[
\|x+y\|^2 - \|x-y\|^2 = 4\ip{x}{y}.
\]

Dividing by 4 gives the result. \qed
\end{mdframed}

\bigskip

%% ============================================================
%% HW 18: Law of cosines in R^2
%% ============================================================

\needspace{4\baselineskip}
\hfill \break  $~\!\!$ \dotfill \medskip \\
\textbf{Question 18.}\\
Show that the identity $\ip{x}{y} = \frac{1}{2}(\|x+y\|^2 - \|x\|^2 - \|y\|^2)$ is the law of cosines in $X = \RR^2$ with the dot product. Also show that $\theta$ (the angle between $x$ and $y$) is the usual geometric angle.

\noindent\\
\textbf{Solution.}\\
\quad\\
\begin{mdframed}[linewidth=1pt]
In $\RR^2$ with the dot product, $\ip{x}{y} = x \cdot y = \|x\|\|y\|\cos\theta$ where $\theta$ is the angle between $x$ and $y$.

From the identity:
\[
\ip{x}{y} = \frac{1}{2}\left(\|x+y\|^2 - \|x\|^2 - \|y\|^2\right).
\]

Now $\|x + y\|^2 = \|x\|^2 + 2(x \cdot y) + \|y\|^2$, so
\[
\frac{1}{2}(\|x+y\|^2 - \|x\|^2 - \|y\|^2) = x \cdot y. \quad \checkmark
\]

For the law of cosines: let $a = \|x\|$, $b = \|y\|$, $c = \|x - y\|$ (the side opposite the angle $\theta$). Then
\[
c^2 = \|x - y\|^2 = a^2 + b^2 - 2\ip{x}{y} = a^2 + b^2 - 2ab\cos\theta,
\]
which is exactly the classical law of cosines. Hence the abstract angle $\theta = \arccos\frac{\ip{x}{y}}{\|x\|\|y\|}$ agrees with the usual geometric angle. \qed
\end{mdframed}

\bigskip

%% ============================================================
%% HW 19: Isometry C case
%% ============================================================

\needspace{4\baselineskip}
\hfill \break  $~\!\!$ \dotfill \medskip \\
\textbf{Question 19.}\\
Prove that if $T: V \to W$ is a linear map between inner product spaces over $\CC$ and $\|T(v)\|_W = \|v\|_V$ for all $v$, then $T$ is an isometry (i.e., $\ip{T(u)}{T(v)}_W = \ip{u}{v}_V$).

(Hint: use the polarization identity for $\CC$.)

\noindent\\
\textbf{Solution.}\\
\quad\\
\begin{mdframed}[linewidth=1pt]
The polarization identity over $\CC$ gives:
\[
\ip{u}{v} = \frac{1}{4}\left(\|u+v\|^2 - \|u-v\|^2 + i\|u+iv\|^2 - i\|u-iv\|^2\right).
\]

Since $T$ is linear and norm-preserving:
\begin{align*}
\ip{T(u)}{T(v)}_W &= \frac{1}{4}\Big(\|T(u)+T(v)\|^2 - \|T(u)-T(v)\|^2 \\
&\qquad + i\|T(u)+iT(v)\|^2 - i\|T(u)-iT(v)\|^2\Big) \\
&= \frac{1}{4}\Big(\|T(u+v)\|^2 - \|T(u-v)\|^2 + i\|T(u+iv)\|^2 - i\|T(u-iv)\|^2\Big) \\
&= \frac{1}{4}\Big(\|u+v\|^2 - \|u-v\|^2 + i\|u+iv\|^2 - i\|u-iv\|^2\Big) \\
&= \ip{u}{v}_V.
\end{align*}

Hence $T$ is an isometry. \qed
\end{mdframed}

\end{document}
